\section{Formalización de la Representación Matricial Relativa}

La piedra angular de la presente investigación reside en la transición de un modelo de datos basado en eventos absolutos —típico de los sistemas de Inteligencia Artificial que operan sobre el estándar MIDI tradicional— hacia un paradigma de representación matricial relativa. Esta formalización no busca meramente codificar la altura de las notas, sino capturar la semántica subyacente de la armonía funcional y la jerarquía tonal. 

Como se observa en el desarrollo del módulo \texttt{procesador\_midi.py}, la música no se trata como una sucesión estadística de números de nota (0-127), sino como un sistema de relaciones vectoriales. La importancia de este enfoque radica en la \textit{invarianza tonal}: la capacidad de un modelo de aprendizaje profundo para reconocer una progresión armónica o una línea melódica independientemente de la tónica en la que se ejecute. Según \cite{tymoczko2011}, la música occidental se organiza en espacios geométricos donde la proximidad no es solo acústica, sino funcional; nuestra representación matricial formaliza esta proximidad mediante el cálculo dinámico de grados y distancias interválicas.

\subsection{Abstracción de datos: Del evento absoluto al vector relativo}

El proceso de abstracción comienza con la deconstrucción del archivo MIDI. Mientras que el estándar MIDI define una nota mediante un entero absoluto $p \in [0, 127]$, nuestro sistema descompone esta variable en dos componentes críticos: la \textbf{Clase de Altura} (\textit{Pitch Class}) y el \textbf{Grado Escalar}.

En el código fuente de \texttt{procesador\_midi.py}, esta abstracción se materializa en la extracción inicial de la clase de altura mediante la operación de módulo:

\begin{verbatim}
# Extracción de la clase de altura en procesador_midi.py
midi_pitch = note['pitch']
pitch_class = midi_pitch % 12
row.append(float(pitch_class))
\end{verbatim}

Sin embargo, el \texttt{pitch\_class} por sí solo es insuficiente para la composición profesional. La abstracción propuesta añade una dimensión de temporalidad relativa. En lugar de utilizar marcas de tiempo absolutas (ticks), el sistema normaliza la duración basándose en la resolución del archivo (\texttt{ticks\_per\_beat}), permitiendo que el modelo aprenda ritmos independientemente del \textit{tempo} (BPM) seleccionado:

\begin{verbatim}
# Normalización rítmica relativa
duration_units = note['duration'] / ticks_per_unit
row.append(round(duration_units, 2))
\end{verbatim}

Esta primera capa de abstracción transforma un evento discreto en un vector bi-dimensional inicial $[PC, D]$, donde $PC$ es la clase de altura y $D$ es la duración relativa. Este paso es fundamental para mitigar el problema de la dispersión de datos en redes neuronales; al reducir el dominio de 128 notas a 12 clases de altura normalizadas rítmicamente, se facilita la convergencia del entrenamiento \cite{briot2020deep}.

\subsection{Algoritmo de mapeo de funciones armónicas y distancias interválicas}

El núcleo innovador del conversor es el algoritmo de mapeo de grados escalares. A diferencia de modelos como \textit{MusicTransformer} \cite{huang2018music}, que dependen de la atención posicional para inferir la tonalidad, nuestro sistema inyecta el conocimiento teórico de la escala directamente en la matriz de entrada.

La función \texttt{get\_scale\_degree} implementa un mapeo matemático que proyecta la nota actual sobre el espacio de la tónica definida. Sea $P$ la clase de altura de la nota, $T$ la tónica y $S$ el conjunto de intervalos de la escala (ej. $S_{Mayor} = \{0, 2, 4, 5, 7, 9, 11\}$). El Grado Escalar $G$ se define como:

\begin{equation}
R = (P - T) \mod 12
\end{equation}
\begin{equation}
G = 
\begin{cases} 
index(R, S) + 1 & \text{si } R \in S \\
(i + 1) + 0.5 & \text{si } S_i < R < S_{i+1} \\
|S| + 0.5 & \text{si } R > max(S)
\end{cases}
\end{equation}

Esta lógica permite manejar no solo notas diatónicas (grados enteros), sino también notas cromáticas o "alteraciones accidentales" (grados fraccionarios), lo cual es esencial para el contrapunto avanzado y la armonía de intercambio modal. El código implementa esta jerarquía de la siguiente manera:

\begin{verbatim}
def get_scale_degree(pitch_class, root_note, scale_intervals):
    relative_pitch = (pitch_class - root_note) % 12
    
    if relative_pitch in scale_intervals:
        return float(scale_intervals.index(relative_pitch) + 1)
    
    for i in range(len(scale_intervals) - 1):
        lower = scale_intervals[i]
        upper = scale_intervals[i+1]
        if lower < relative_pitch < upper:
            return float((i + 1) + 0.5)
    # ... (manejo de bordes)
\end{verbatim}

Adicionalmente, se integra el cálculo de la \textbf{distancia interválica relativa} entre notas consecutivas. Esta columna de la matriz permite que el algoritmo genético evalúe la "suavidad" de la conducción de voces (\textit{voice leading}). En \texttt{procesador\_midi.py}, el intervalo se normaliza a unidades de tono (donde 0.5 es un semitono), alineándose con la teoría de \cite{chew2014mathematical} sobre el modelado matemático de la tensión melódica:

\begin{verbatim}
# Cálculo de intervalo relativo en tonos
semitones = midi_pitch - last_note_pitch
interval = semitones / 2.0
row.append(interval)
\end{verbatim}

\subsection{Normalización y preparación del dataset para el entrenamiento}

La fase final de la formalización es la estructuración de la matriz resultante para su consumo por arquitecturas de aprendizaje profundo (RNN/LSTM). La matriz final posee una topología de $N \times 4$, donde cada fila representa un evento musical con cuatro parámetros paramétricos:
\begin{enumerate}
    \item \textbf{Clase de Altura (Absoluta Octava):} Proporciona el color de la nota dentro del total cromático.
    \item \textbf{Duración (Relativa):} Cuantiza el ritmo en unidades de tiempo musical (corcheas, semicorcheas).
    \item \textbf{Grado Escalar (Funcional):} Define la posición de la nota respecto a la jerarquía de la escala, facilitando el aprendizaje de resoluciones armónicas.
    \item \textbf{Intervalo (Melódico):} Captura el movimiento horizontal de la voz.
\end{enumerate}

La preparación del dataset implica la vectorización de estas matrices. Dado que las redes recurrentes requieren secuencias de longitud fija, se implementa un sistema de \textit{sliding window} (ventana deslizante). Sin embargo, a diferencia de otros datasets simbólicos que solo guardan la nota, nuestra preparación conserva la \textbf{metadata de la tonalidad}. 

Esto permite realizar una \textit{normalización tonal activa}: durante el entrenamiento, se pueden trasponer todas las piezas a una tónica común de forma matemática, garantizando que el modelo aprenda que un grado "5.0" (Dominante) siempre tiende a resolver a un "1.0" (Tónica), independientemente de si la pieza original estaba en Fa sostenido o Do mayor. Este proceso es fundamental para alcanzar el objetivo de "generación paramétrica estructuralmente coherente" planteado en el Capítulo 1, ya que reduce drásticamente el espacio de búsqueda del algoritmo genético al operar sobre un dataset pre-estructurado bajo leyes armónicas \cite{hadjeres2017deepbach}.

\begin{verbatim}
# Estructura final de la fila de la matriz para entrenamiento
# [PitchClass, DurationUnits, ScaleDegree, IntervalTones]
result_matrix = midi_to_matrix(target_file, root_note_name='C', scale_type='major')
\end{verbatim}

En conclusión, el conversor implementado no es solo una herramienta de transformación de formato, sino un motor de extracción de características teóricas. La representación resultante permite que el sistema híbrido de IA actúe como un compositor instruido que "entiende" la función de cada nota dentro de un contexto tonal, superando las limitaciones de la generación estocástica convencional.
